\documentclass[11pt,a4paper]{article}

\usepackage[margin=2.5cm]{geometry}
\usepackage{amsmath}
\usepackage{amssymb}
\usepackage{bm}
\usepackage{hyperref}
\usepackage[ruled,vlined]{algorithm2e}

\newcommand{\W}{\mathbf{W}}
\newcommand{\Wt}{\widetilde{\mathbf{W}}}
\newcommand{\xx}{\mathbf{x}}
\newcommand{\y}{\mathbf{y}}
\newcommand{\yt}{\widetilde{\mathbf{y}}}
\newcommand{\beps}{\boldsymbol{\epsilon}}
\newcommand{\bgam}{\boldsymbol{\gamma}}
\newcommand{\bbet}{\boldsymbol{\beta}}
\newcommand{\E}{\mathbb{E}}
\newcommand{\Var}{\mathrm{Var}}
\newcommand{\norm}[1]{\mathcal{N}\!\left(#1\right)}

\title{Analytical Bias and Variance Correction\\for Quantized Linear and Convolutional Layers}
\author{ACIQ project (derivation, after Nagel et al.\ 2019)}
\date{}

\begin{document}
\maketitle

\section{Setup and Notation}

We consider a single weight-bearing layer that produces pre-activation output $\y$ from input $\xx$ using FP32 weights $\W$ and bias $\mathbf{b}$. After post-training quantization the same layer uses quantized weights $\Wt$. Bias is kept in FP32. We define the per-element \emph{quantization error}
\begin{equation}
\beps \;=\; \Wt - \W,
\end{equation}
so that
\begin{equation}
\yt \;=\; \Wt\,\xx + \mathbf{b} \;=\; \W\,\xx + \beps\,\xx + \mathbf{b} \;=\; \y + \beps\,\xx.
\label{eq:basic_decomposition}
\end{equation}

We assume Batch Normalization (BN) is folded into the preceding weight layer, but that the BN scale~$\bgam$ and shift~$\bbet$ have been captured before fusion. The activation function is ReLU. Under the assumption made in~\cite{nagel2019} that the pre-ReLU output of a BN-folded layer is approximately Gaussian, the per-channel input distribution to a downstream layer follows a \emph{ReLU-clipped normal} with parameters $(\bbet_c, \bgam_c^2)$.

Throughout this document we use:
\begin{itemize}
\item $c_o$, $c_i$ for output / input channel indices.
\item $m, n$ for kernel spatial indices.
\item $\E[\xx_{c_i}]$ and $\Var[\xx_{c_i}]$ for the per-channel mean and variance of the layer input, assumed spatially stationary (i.i.d.\ across spatial positions within a channel).
\item Independence \emph{between} input channels (a standard analytical simplification; the same as~\cite{nagel2019}).
\end{itemize}

\section{Bias Correction --- Linear Layer}
\label{sec:bias_linear}

For a fully-connected layer~\cite[Eq.~16--17]{nagel2019},
\begin{align}
\E[\y] \;&=\; \E[\y] + \E[\beps\,\xx] - \E[\beps\,\xx] \\
       \;&=\; \E[\yt] - \E[\beps\,\xx],
\end{align}
so subtracting $\E[\beps\,\xx] = \beps\,\E[\xx]$ from $\yt$ recovers the FP32 mean. Per output unit $c_o$,
\begin{equation}
\Delta b_{c_o} \;=\; \sum_{c_i} \beps_{c_o,c_i}\,\E[\xx_{c_i}].
\end{equation}
The correction is folded into the bias:
\begin{equation}
\mathbf{b}_{c_o}^{\,\mathrm{corr}} \;=\; \mathbf{b}_{c_o} - \Delta b_{c_o}.
\label{eq:bias_correction_linear}
\end{equation}

\section{Bias Correction --- Convolutional Layer with ReLU}
\label{sec:bias_conv}

For a 2D convolution $\y_{c_o,i,j} = \sum_{c_i,m,n} \W_{c_o,c_i,m,n}\,\xx_{c_i,i-m,j-n} + \mathbf{b}_{c_o}$, the same decomposition gives $\E[\beps\!*\!\xx]$ pointwise. Under spatial stationarity $\E[\xx_{c_i,i-m,j-n}] = \E[\xx_{c_i}]$:
\begin{align}
\bigl[\beps\!*\!\E[\xx]\bigr]_{c_o,i,j}
  &= \sum_{c_i,m,n}\E[\xx_{c_i,i-m,j-n}]\,\beps_{c_o,c_i,m,n} \\
  &= \sum_{c_i}\!\left[\E[\xx_{c_i}]\,\sum_{m,n}\beps_{c_o,c_i,m,n}\right].
\label{eq:bias_correction_conv_expanded}
\end{align}
This is independent of the spatial location $(i,j)$ and can be folded into the per-output-channel bias. In vector form, with $\beps_\Sigma \in \mathbb{R}^{C_o\times C_i}$ defined by $[\beps_\Sigma]_{c_o,c_i} = \sum_{m,n}\beps_{c_o,c_i,m,n}$,
\begin{equation}
\boxed{\;\Delta\mathbf{b} \;=\; \beps_\Sigma\,\E[\xx], \qquad \mathbf{b}^{\,\mathrm{corr}} \;=\; \mathbf{b} - \Delta\mathbf{b}.\;}
\label{eq:bias_correction_conv_final}
\end{equation}
Since~\eqref{eq:bias_correction_linear} is the special case $K=1$ of \eqref{eq:bias_correction_conv_final} (or equivalently $\beps_\Sigma = \beps$ for a fully-connected layer), the implementation handles both with a single contraction.

\section{Variance Correction --- Linear Layer}
\label{sec:var_linear}

The paper of Nagel et al.\ corrects only the mean. However, weight quantization perturbs the per-channel variance as well. Under per-channel input independence,
\begin{equation}
\Var[\y_{c_o}] \;=\; \sum_{c_i}\W_{c_o,c_i}^2\,\Var[\xx_{c_i}], \qquad
\Var[\yt_{c_o}] \;=\; \sum_{c_i}\Wt_{c_o,c_i}^2\,\Var[\xx_{c_i}].
\end{equation}
Define the per-output-channel rescaling
\begin{equation}
s_{c_o} \;=\; \sqrt{\,\Var[\y_{c_o}]\,/\,\Var[\yt_{c_o}]\,}.
\end{equation}
We apply a \emph{mean-preserving} affine rescaling so that the second moment matches FP32 without disturbing whatever mean is in place:
\begin{equation}
\yt_{c_o}^{\,\mathrm{corr}} \;=\; s_{c_o}\bigl(\yt_{c_o} - \E[\yt_{c_o}]\bigr) + \E[\yt_{c_o}]
\;=\; s_{c_o}\,\yt_{c_o} + (1-s_{c_o})\,\E[\yt_{c_o}].
\label{eq:var_correction_form}
\end{equation}
Folding $s_{c_o}$ into the per-channel weight row and into the bias gives:
\begin{equation}
\boxed{\;\Wt^{\,\mathrm{corr}}_{c_o} \;=\; s_{c_o}\,\Wt_{c_o}, \qquad
\mathbf{b}^{\,\mathrm{corr}}_{c_o} \;=\; s_{c_o}\,\mathbf{b}_{c_o} + (1-s_{c_o})\,\E[\yt_{c_o}].\;}
\label{eq:var_correction_linear}
\end{equation}
Here $\E[\yt_{c_o}] = \sum_{c_i}\Wt_{c_o,c_i}\,\E[\xx_{c_i}] + \mathbf{b}_{c_o}$ is the (uncorrected) quantized mean.

\section{Variance Correction --- Convolutional Layer with ReLU}

Let $\Wt^{(2)}_\Sigma \in \mathbb{R}^{C_o\times C_i}$ with $[\Wt^{(2)}_\Sigma]_{c_o,c_i} = \sum_{m,n}\Wt_{c_o,c_i,m,n}^2$, and similarly for $\W^{(2)}_\Sigma$. Under the same independence and spatial-i.i.d.\ assumptions used by~\cite[App.~B]{nagel2019},
\begin{equation}
\Var[\y_{c_o,i,j}] \;=\; \sum_{c_i}[\W^{(2)}_\Sigma]_{c_o,c_i}\,\Var[\xx_{c_i}],
\end{equation}
which is independent of the spatial location $(i,j)$. Defining $s_{c_o}$ as before, the same folding as~\eqref{eq:var_correction_linear} holds, with the convention that $s_{c_o}$ broadcasts across the kernel spatial axes:
\begin{equation}
\boxed{\;\Wt^{\,\mathrm{corr}}_{c_o,c_i,m,n} \;=\; s_{c_o}\,\Wt_{c_o,c_i,m,n}, \qquad
\mathbf{b}^{\,\mathrm{corr}}_{c_o} \;=\; s_{c_o}\,\mathbf{b}_{c_o} + (1-s_{c_o})\,\E[\yt_{c_o}],\;}
\end{equation}
with $\E[\yt_{c_o}] = \sum_{c_i}[\Wt_\Sigma]_{c_o,c_i}\,\E[\xx_{c_i}] + \mathbf{b}_{c_o}$.

\paragraph{Numerical guard.} If $\Var[\yt_{c_o}]$ is near-zero (a near-dead channel), $s_{c_o}$ explodes. We clip $s_{c_o}\in[0.5,2.0]$ in the implementation.

\section{Joint Bias \& Variance Correction}
\label{sec:joint}

When both corrections are applied, the variance scaling must pivot around the bias-corrected (FP32) mean rather than around the uncorrected quantized mean. Let
\begin{equation}
\mathbf{b}^{\,\mathrm{bias}} \;=\; \mathbf{b} - \Delta\mathbf{b}
\quad\text{and}\quad
\E[\y_{c_o}] \;=\; \sum_{c_i}[\W_\Sigma]_{c_o,c_i}\,\E[\xx_{c_i}] + \mathbf{b}_{c_o},
\end{equation}
where $\Delta\mathbf{b}$ is from~\eqref{eq:bias_correction_conv_final} and $\E[\y_{c_o}]$ is the FP32 target mean. The variance scaling is then applied around $\E[\y_{c_o}]$:
\begin{equation}
\boxed{\;\Wt^{\,\mathrm{joint}}_{c_o} \;=\; s_{c_o}\,\Wt_{c_o}, \qquad
\mathbf{b}^{\,\mathrm{joint}}_{c_o} \;=\; s_{c_o}\bigl(\mathbf{b}_{c_o} - \Delta b_{c_o}\bigr) + (1-s_{c_o})\,\E[\y_{c_o}].\;}
\label{eq:joint_correction}
\end{equation}
A direct check confirms that $\E[\y^{\,\mathrm{joint}}_{c_o}] = \E[\y_{c_o}]$ and $\Var[\y^{\,\mathrm{joint}}_{c_o}] = \Var[\y_{c_o}]$ under the working assumptions.

\section{Computing $\E[\xx]$ and $\Var[\xx]$ from Batch-Norm Parameters}
\label{sec:clipped_normal}

When BN is applied before ReLU, the pre-ReLU output of channel $c$ is approximately $\norm{\mu=\bbet_c,\,\sigma^2=\bgam_c^2}$. The post-ReLU input to the next layer is then the \emph{clipped normal} on $[a,b]=[0,\infty)$. With $\alpha = -\bbet_c/\bgam_c$:
\begin{align}
\E[\xx_c] \;&=\; \bgam_c\,\varphi(\alpha) + \bbet_c\bigl(1 - \Phi(\alpha)\bigr), \label{eq:Ex_relu}\\
\Var[\xx_c] \;&=\; (\bbet_c^2 + \bgam_c^2)\bigl(1-\Phi(\alpha)\bigr) + \bgam_c\bbet_c\,\varphi(\alpha) - \E[\xx_c]^2,
\label{eq:Varx_relu}
\end{align}
where $\varphi(\cdot)$ is the standard normal PDF and $\Phi(\cdot)$ the standard normal CDF~\cite[App.~C]{nagel2019}.

For the general clipped normal $f(X) = \mathrm{clip}(X, a, b)$ with $X \sim \norm{\mu,\sigma^2}$ and $\alpha = (a-\mu)/\sigma$, $\beta_n=(b-\mu)/\sigma$, $Z = \Phi(\beta_n)-\Phi(\alpha)$, the closed-form mean is
\begin{equation}
\E[f(X)] = \sigma\bigl(\varphi(\alpha)-\varphi(\beta_n)\bigr) + \mu Z + a\Phi(\alpha) + b\bigl(1 - \Phi(\beta_n)\bigr),
\end{equation}
and the variance follows from the second-moment formula in~\cite[App.~C]{nagel2019}. The implementation handles the general case with $a=0,\,b=\infty$ as the ReLU specialization.

\section{Propagating Statistics Through Residual Connections and GAP}
\label{sec:residual}

Within a ResNet block, $\E[\xx]$ and $\Var[\xx]$ propagate as follows.

\paragraph{In-block conv inputs.} For \texttt{conv2} (and \texttt{conv3} in a Bottleneck), the input is $\mathrm{ReLU}(\mathrm{BN}_{\text{prev conv}}(\cdot))$, so~\eqref{eq:Ex_relu}--\eqref{eq:Varx_relu} apply with $(\bbet,\bgam)$ taken from the immediately preceding BN.

\paragraph{Residual sum.} The output of the block is $\mathrm{ReLU}\bigl(\mathrm{BN}_{\text{main}}(\cdot) + \text{skip}\bigr)$. Treating $\text{skip}$ as approximately Gaussian and independent of the main branch (a simplification used by~\cite{nagel2019}), the pre-ReLU per-channel distribution is
\begin{equation}
\norm{\bbet_{\text{main}} + \mu_{\text{skip}},\;\bgam_{\text{main}}^2 + \sigma_{\text{skip}}^2},
\end{equation}
and the post-ReLU activation statistics follow from~\eqref{eq:Ex_relu}--\eqref{eq:Varx_relu}.
\begin{itemize}
\item \emph{Identity skip:} $(\mu_{\text{skip}}, \sigma_{\text{skip}}^2)$ is the previous block's post-residual stats (recursive).
\item \emph{Downsample skip:} the skip path is also $\mathrm{BN}_{\text{ds}}$, contributing $(\bbet_{\text{ds}},\bgam_{\text{ds}}^2)$.
\end{itemize}

\paragraph{Global average pooling (\texttt{fc} input).} The fully-connected layer of a ResNet sees $\bar{\xx}_c = \frac{1}{HW}\sum_{i,j}\xx_{c,i,j}$. Under the spatial-i.i.d.\ assumption,
\begin{equation}
\E[\bar{\xx}_c] \;=\; \E[\xx_c], \qquad \Var[\bar{\xx}_c] \;=\; \Var[\xx_c]\,/\,(HW),
\end{equation}
so the FC layer uses the layer-4 output mean unchanged and a variance scaled by $1/(HW)$ (e.g.\ $1/49$ for a $7\times 7$ feature map at $224\times 224$ input).

\paragraph{Stem.} The first convolution receives the normalized image, not the output of a BN/ReLU. Analytical correction is undefined here and the stem layer is left uncorrected in the implementation.

\section{Algorithm}

\begin{algorithm}[H]
\SetAlgoLined
\KwIn{FP32 model with BN; quantization method; correction mode $\in\{\text{none, bias, variance, joint}\}$.}
\KwOut{Corrected quantized model.}
$\{(\bgam_b,\bbet_b)\}_b \leftarrow$ \texttt{capture\_bn\_params(model)}\;
fold BN into Conv (in place)\;
$\{(\E[\xx],\Var[\xx])\}_L \leftarrow$ \texttt{compute\_input\_stats}, propagating per~\S\ref{sec:residual}\;
\For{each weight layer $L$ except stem}{
  Save $\W^{\,\mathrm{fp}}_L$\;
  Quantize: $\Wt_L$\;
  $\beps_\Sigma \leftarrow (\Wt - \W^{\,\mathrm{fp}}).\mathrm{sum}(\text{kernel axes})$\;
  $\Delta\mathbf{b} \leftarrow \beps_\Sigma\cdot\E[\xx]_L$\;
  $\mathbf{s} \leftarrow \sqrt{\,(\W^{\,\mathrm{fp}\,(2)}_\Sigma\,\Var[\xx])\,/\,(\Wt^{(2)}_\Sigma\,\Var[\xx])\,}$, clipped\;
  \uIf{mode = \emph{none}}{no change}
  \uElseIf{mode = \emph{bias}}{$\mathbf{b} \leftarrow \mathbf{b} - \Delta\mathbf{b}$}
  \uElseIf{mode = \emph{variance}}{
    $\E[\yt] \leftarrow \Wt_\Sigma\cdot\E[\xx] + \mathbf{b}$\;
    $\Wt \leftarrow \mathbf{s}\odot\Wt$;\; $\mathbf{b} \leftarrow \mathbf{s}\odot\mathbf{b} + (\mathbf{1}-\mathbf{s})\odot\E[\yt]$\;
  }
  \Else(\emph{joint}){
    $\E[\y] \leftarrow \W^{\,\mathrm{fp}}_\Sigma\cdot\E[\xx] + \mathbf{b}$\;
    $\Wt \leftarrow \mathbf{s}\odot\Wt$;\; $\mathbf{b} \leftarrow \mathbf{s}\odot(\mathbf{b}-\Delta\mathbf{b}) + (\mathbf{1}-\mathbf{s})\odot\E[\y]$\;
  }
}
\caption{Analytical bias / variance correction for a quantized model.}
\end{algorithm}

\section{Implementation Notes}

\begin{itemize}
\item \textbf{Variance and joint correction require per-channel weight quantization.} The scale $s_{c_o}$ is per output channel by construction. Per-channel weight quantization stores one scale per output channel, so $s_{c_o}$ folds cleanly into that scale (or, equivalently in this dequantized-FP32 pipeline, into the per-channel weight row). Per-tensor weight quantization stores a single shared scale; multiplying by a per-channel $s_{c_o}$ produces a per-channel grid that is no longer representable under per-tensor quant. To keep the benchmark over hardware-realizable variants only, the implementation restricts $\{$bias, variance, joint$\}$ corrections to per-channel quantization variants. Per-tensor variants run only as the uncorrected baseline. (A hardware-equivalent alternative would be to express $s_{c_o}$ as a per-output-channel rescale at the accumulator/requant boundary, leaving the integer weight tensor untouched; that is supported by typical INT8 inference engines but is out of scope here.)
\item Bias correction is hardware-friendly in any quantization scheme: it is an additive shift on the bias tensor, which is FP32/INT32 in any standard pipeline.
\item The independence assumption between input channels is the same simplification used by~\cite{nagel2019} for $\E[\xx]$. It is the dominant approximation; real activations exhibit modest cross-channel correlation, which would add cross terms to $\Var[\y_{c_o}]$ that this derivation omits.
\item The analytical clipped-normal model uses the BN running statistics. After fusion, the BN parameters $(\bgam,\bbet)$ are captured beforehand; the running mean and variance are absorbed into the fused weight and bias.
\item The stem (first) convolution receives the normalized image and is left uncorrected: there is no preceding BN/ReLU to give an analytical input distribution.
\end{itemize}

\begin{thebibliography}{1}
\bibitem{nagel2019}
Markus Nagel, Mart van Baalen, Tijmen Blankevoort, Max Welling.
\emph{Data-Free Quantization Through Weight Equalization and Bias Correction.}
ICCV 2019. \url{https://arxiv.org/abs/1906.04721}
\end{thebibliography}

\end{document}
